\documentclass[11pt,a4paper,UTF8]{ctexart}
\usepackage{geometry}
\geometry{left=18mm,right=18mm,top=24mm,bottom=24mm,headheight=22pt,footskip=12mm}
\usepackage{amsmath,amssymb,mathtools,bm,esint,extarrows}
\usepackage{xcolor}
\usepackage{tcolorbox}
\tcbuselibrary{skins,breakable}
\usepackage{fancyhdr}
\usepackage{titlesec}
\usepackage{hyperref}
\usepackage{tikz}
\usepackage{booktabs}
\usepackage{tabularx}

% --- Color Palette (Purple & DeepPink Academic Theme) ---
\definecolor{DeepPurple}{HTML}{4C1D95}
\definecolor{TitlePurple}{HTML}{581C87}
\definecolor{SubPurple}{HTML}{6B21A8}
\definecolor{DeepPink}{HTML}{FF1493}
\definecolor{LightPinkBg}{HTML}{FFF5F8}
\definecolor{LightPurpleBg}{HTML}{F8F5FF}
\definecolor{GoldAccent}{HTML}{D97706}
\definecolor{DarkText}{HTML}{1F2937}

\hypersetup{
    colorlinks=true,
    linkcolor=TitlePurple,
    citecolor=DeepPink,
    urlcolor=SubPurple,
    pdfborder={0 0 0}
}

% --- Typography & Spacing ---
\linespread{1.3}
\setlength{\parskip}{0.35em plus 0.1em minus 0.1em}
\setlength{\parindent}{0pt}

% --- Section Titles Styling ---
\titleformat{\section}
  {\Large\bfseries\color{TitlePurple}}{\thesection}{1em}{}
  [\vspace{1mm}{\color{TitlePurple!30}\hrule height 1pt}]
\titleformat{\subsection}
  {\large\bfseries\color{SubPurple}}{\thesubsection}{0.8em}{}
\titleformat{\subsubsection}
  {\normalsize\bfseries\color{DeepPurple}}{\thesubsubsection}{0.6em}{}

% --- tcolorbox Environments ---
% 1. Theorem / Formula / Method / Analysis Box (Deep Purple)
\newtcolorbox{thmbox}[1][]{
    enhanced,
    breakable,
    colback=LightPurpleBg,
    colframe=TitlePurple,
    arc=3pt,
    boxrule=0pt,
    leftrule=3.5pt,
    left=10pt,
    right=10pt,
    top=8pt,
    bottom=8pt,
    title=#1,
    coltitle=white,
    colbacktitle=TitlePurple,
    fonttitle=\bfseries\small,
    attach boxed title to top left={yshift=-2mm, xshift=4mm},
    boxed title style={boxrule=0pt, arc=2pt}
}

% 2. Solution / Formula / Derivation Box (DeepPink)
\newtcolorbox{solbox}[1][]{
    enhanced,
    breakable,
    colback=LightPinkBg,
    colframe=DeepPink,
    arc=3pt,
    boxrule=0pt,
    leftrule=3.5pt,
    left=10pt,
    right=10pt,
    top=8pt,
    bottom=8pt,
    title=#1,
    coltitle=white,
    colbacktitle=DeepPink,
    fonttitle=\bfseries\small,
    attach boxed title to top left={yshift=-2mm, xshift=4mm},
    boxed title style={boxrule=0pt, arc=2pt}
}

% 3. Learning Goals / Summary Box
\newtcolorbox{targetbox}[1][]{
    enhanced,
    breakable,
    colback=LightPurpleBg,
    colframe=DeepPurple,
    arc=3pt,
    boxrule=1pt,
    left=10pt,
    right=10pt,
    top=8pt,
    bottom=8pt,
    title=#1,
    coltitle=white,
    colbacktitle=DeepPurple,
    fonttitle=\bfseries\small,
    attach boxed title to top left={yshift=-2mm, xshift=4mm},
    boxed title style={boxrule=0pt, arc=2pt}
}

\begin{document}

% ------------------- 讲义封面 (Cover Page) -------------------
\begin{titlepage}
\thispagestyle{empty}
\begin{tikzpicture}[remember picture,overlay]
    \fill[DeepPurple] (current page.north west) rectangle ([yshift=-7cm]current page.north east);
    \draw[DeepPink, line width=3.5pt] ([yshift=-7cm]current page.north west) -- ([yshift=-7cm]current page.north east);
    \draw[GoldAccent, line width=1pt] ([yshift=-7.12cm]current page.north west) -- ([yshift=-7.12cm]current page.north east);
    
    \node[anchor=north east, opacity=0.12, text=white] at ([xshift=-1.5cm, yshift=-0.5cm]current page.north east) {
        \fontsize{70}{70}\selectfont\bfseries 2027
    };
\end{tikzpicture}

\vspace*{0.8cm}
\begin{center}
    {\color{white}\fontsize{26}{32}\selectfont\bfseries 2027 考研数学(二) 核心讲义}\\[0.6cm]
    {\color{white}\fontsize{13}{18}\selectfont\bfseries 全国硕士研究生招生考试 · 统考数学(二) 核心精讲全书}\\[1.5cm]
\end{center}

\vspace{3.5cm}
\begin{center}
    {\color{GoldAccent}\large\bfseries $\blacktriangleright$ 基础强化 · 核心题解 · 考点深水区 $\blacktriangleleft$}\\[0.8cm]
    {\color{TitlePurple}\fontsize{24}{28}\selectfont\bfseries [强化]一元积分学的计算}\\[0.6cm]
    {\color{DeepPink}\large\bfseries —— 考研高阶冲刺与深水区专攻 ——}\\[1.5cm]
\end{center}

\vfill

\begin{center}
\begin{tcolorbox}[
    colback=white,
    colframe=DeepPurple,
    width=0.88\textwidth,
    arc=4pt,
    boxrule=1.2pt,
    halign=center,
    drop shadow=gray!30
]
    \vspace{3mm}
    {\bfseries\large 编著团队：EscoffierZhou}\\[2.5mm]
    {\bfseries\color{TitlePurple} 目标院校：中国科学院大学 (UCAS) 计算机专硕 (22408)}\\[2.5mm]
    {\small\color{gray} 备考铁律：手算真题数字 · 错题白纸重做 · 408 客观题为王}\\[2.5mm]
    {\small 编制版本：2027 届首发版 · \today}
    \vspace{2mm}
\end{tcolorbox}
\end{center}
\vspace{0.8cm}
\end{titlepage}

% ------------------- 目录与页面主体配置 -------------------
\pagestyle{fancy}
\fancyhf{}
\fancyhead[L]{\small\color{gray}2027 考研数学(二) · 核心讲义系列}
\fancyhead[R]{\small\color{TitlePurple}\nouppercase{\leftmark}}
\fancyfoot[C]{\small\color{gray}— 第 \thepage\ 页 —}
\fancyfoot[R]{\small\color{gray}UCAS 22408}
\renewcommand{\headrulewidth}{0.5pt}

\pagenumbering{Roman}
\tableofcontents
\newpage
\pagenumbering{arabic}


\begin{targetbox}[本章复习目标与核心知识图谱]
\begin{itemize}
  \item 目的[1]:掌握欧拉代换(三大代换模型)与配对积分法，秒杀复杂二次根式与三角有理式
\end{itemize}
\end{targetbox}

\begin{targetbox}[本章复习目标与核心知识图谱]
\begin{itemize}
  \item 目的[2]:深刻理解广义区间再现公式与柯西-阿罗积分模型，攻克高难度定积分计算
\end{itemize}
\end{targetbox}

\begin{targetbox}[本章复习目标与核心知识图谱]
\begin{itemize}
  \item 目的[3]:掌握柯西-施瓦茨(Cauchy-Schwarz)积分不等式与变限积分微分方程转化法
\end{itemize}
\end{targetbox}

\begin{targetbox}[本章复习目标与核心知识图谱]
\begin{itemize}
  \item 目的[4]:理解黎曼-勒贝格引理与高频振荡积分的平均值极限转化规律
\end{itemize}
\end{targetbox}

\begin{targetbox}[本章复习目标与核心知识图谱]
\begin{itemize}
  \item 目的[5]:攻克变限积分内不可积函数的泰勒级数展开内插法与隐函数高阶求导
\end{itemize}
\end{targetbox}

\begin{targetbox}[本章复习目标与核心知识图谱]
\begin{itemize}
  \item 目的[6]:掌握费曼积分号下求导法(参数积分法)与含参反常积分的阶数双侧压制
\end{itemize}
\end{targetbox}


\section{1.不定积分特种深水区代换与秒杀模型}

\begin{thmbox}[模型[1]: 三大欧拉代换(Euler Substitutions)终极消根]

欧拉代换是解决形如 $\int R(x, \sqrt{ax^2+bx+c})\,\mathrm{d}x$ 的终极通用武器，尤其当被积函数无法通过标准三角代换化简时展现出极强的代数消根威力。

\textbf{[1]} 第一欧拉代换(二次项系数 $a > 0$)：
结构特征: 根号内二次项系数为正，令根式等于 $\pm\sqrt{a}x + t$。


\begin{equation*}
\sqrt{ax^2+bx+c} = \sqrt{a}x + t
\end{equation*}


两边平方：


\begin{equation*}
ax^2 + bx + c = ax^2 + 2\sqrt{a}xt + t^2 \implies bx + c = 2\sqrt{a}xt + t^2
\end{equation*}


二次项 $ax^2$ 完美相消，直接解出一元有理式 $x$：


\begin{equation*}
x = \frac{t^2 - c}{b - 2\sqrt{a}t}
\end{equation*}


求导后 $\mathrm{d}x$ 为有理分式，原积分完全转化为普通的代数有理函数积分。

\textbf{[2]} 第二欧拉代换(常数项系数 $c > 0$)：
结构特征: 根号内常数项为正，令根式等于 $xt \pm \sqrt{c}$。


\begin{equation*}
\sqrt{ax^2+bx+c} = xt + \sqrt{c}
\end{equation*}


两边平方：


\begin{equation*}
ax^2 + bx + c = x^2 t^2 + 2\sqrt{c}xt + c \implies ax^2 + bx = x^2 t^2 + 2\sqrt{c}xt
\end{equation*}


两边约去因子 $x$：


\begin{equation*}
ax + b = xt^2 + 2\sqrt{c}t \implies x = \frac{2\sqrt{c}t - b}{a - t^2}
\end{equation*}


成功将无理根式转化为关于 $t$ 的纯有理分式。

\textbf{[3]} 第三欧拉代换(二次多项式有两个相异实根 $x_1, x_2$)：
结构特征: $ax^2+bx+c = a(x-x_1)(x-x_2)$，令根式等于 $t(x-x_1)$。


\begin{equation*}
\sqrt{a(x-x_1)(x-x_2)} = t(x-x_1)
\end{equation*}


两边平方：


\begin{equation*}
a(x-x_1)(x-x_2) = t^2(x-x_1)^2 \implies a(x-x_2) = t^2(x-x_1)
\end{equation*}


解得：


\begin{equation*}
x = \frac{ax_2 - t^2 x_1}{a - t^2}
\end{equation*}


代回根式即可快速将式子有理化。

\end{thmbox}
\begin{thmbox}[模型[2]: 配对积分法(待定系数线性组合秒杀)]

\textbf{[1]} 正余弦分式线性组合型：
对于积分 $\int \frac{a\sin x + b\cos x}{c\sin x + d\cos x}\,\mathrm{d}x$：

传统万能代换计算量极大，极易出错。高效解法为配对积分法：

设分子为分母与其导数的线性组合：


\begin{equation*}
a\sin x + b\cos x = A(c\sin x + d\cos x) + B(c\sin x + d\cos x)'
\end{equation*}


即：


\begin{equation*}
a\sin x + b\cos x = A(c\sin x + d\cos x) + B(c\cos x - d\sin x)
\end{equation*}


通过比较两端 $\sin x$ 与 $\cos x$ 的系数，联立二元一次方程组求得 $A$ 和 $B$：


\begin{align*}
\begin{cases} cA - dB = a \\ \\ dA + cB = b \end{cases}
\end{align*}


积分瞬间拆为平庸项与对数项：


\begin{equation*}
\int \frac{a\sin x + b\cos x}{c\sin x + d\cos x}\,\mathrm{d}x = A\int 1\,\mathrm{d}x + B\int \frac{\mathrm{d}(c\sin x + d\cos x)}{c\sin x + d\cos x} = Ax + B\ln|c\sin x + d\cos x| + C
\end{equation*}


\textbf{[2]} 指数双曲分式配对型：
对于 $\int \frac{a e^x + b e^{-x}}{c e^x + d e^{-x}}\,\mathrm{d}x$，同理设：


\begin{equation*}
a e^x + b e^{-x} = A(c e^x + d e^{-x}) + B(c e^x + d e^{-x})'
\end{equation*}


可完全仿照上述方法秒解。

\end{thmbox}
\begin{thmbox}[模型[3]: 倒代换高阶复合与双曲代换速算]

\textbf{[1]} 分母高次差幂倒代换：
形如 $\int \frac{\mathrm{d}x}{x^n(1+x^n)^{\frac{1}{n}}}$：

令 $x = \frac{1}{t}$，$\mathrm{d}x = -\frac{1}{t^2}\,\mathrm{d}t$：


\begin{equation*}
= -\int \frac{\frac{1}{t^2}\,\mathrm{d}t}{\frac{1}{t^n}\left(1+\frac{1}{t^n}\right)^{\frac{1}{n}}} = -\int \frac{t^{n-2}\,\mathrm{d}t}{\frac{1}{t}(t^n+1)^{\frac{1}{n}}} = -\int \frac{t^{n-1}\,\mathrm{d}t}{(t^n+1)^{\frac{1}{n}}}
\end{equation*}


分子 $t^{n-1}\,\mathrm{d}t$ 恰好为 $\frac{1}{n}\,\mathrm{d}(t^n+1)$，直接凑微分积出结果。

\textbf{[2]} 双曲代换简化无理积分：
遇到 $\sqrt{x^2+a^2}$，令 $x = a\sinh t$，则 $\sqrt{x^2+a^2} = a\cosh t$，$\mathrm{d}x = a\cosh t\,\mathrm{d}t$。利用 $\cosh^2 t - \sinh^2 t = 1$ 的优良性质，免去三角正割代换后复杂的对数还原。

\end{thmbox}

\section{2.定积分高阶构造、区间再现与深水区不等式}

\begin{thmbox}[模型[1]: 广义区间再现与高难度对称抵消构造]

\textbf{[1]} 中心对称函数的定积分加和定理：
若函数 $f(x)$ 的图像关于点 $(x_0, y_0)$ 成中心对称，即满足 $f(x) + f(2x_0 - x) = 2y_0$。

则在关于 $x_0$ 对称的区间 $[x_0 - a, x_0 + a]$ 上：


\begin{equation*}
\int_{x_0 - a}^{x_0 + a} f(x)\,\mathrm{d}x = 2y_0 a
\end{equation*}


几何解释: 中心对称图形超出水平线 $y = y_0$ 的面积与凹下的面积完全相等，定积分值严格等于矩形面积 $y_0 \cdot 2a$。

\textbf{[2]} 柯西-阿罗经典对数正弦积分模型：

\begin{equation*}
I = \int_0^{\frac{\pi}{2}} \ln(\sin x)\,\mathrm{d}x = \int_0^{\frac{\pi}{2}} \ln(\cos x)\,\mathrm{d}x = -\frac{\pi}{2}\ln 2
\end{equation*}


推导过程:

利用区间再现 $I = \int_0^{\frac{\pi}{2}} \ln(\cos x)\,\mathrm{d}x$。

两式相加：


\begin{equation*}
2I = \int_0^{\frac{\pi}{2}} \ln(\sin x\cos x)\,\mathrm{d}x = \int_0^{\frac{\pi}{2}} \ln\left(\frac{\sin 2x}{2}\right)\,\mathrm{d}x = \int_0^{\frac{\pi}{2}} \ln(\sin 2x)\,\mathrm{d}x - \frac{\pi}{2}\ln 2
\end{equation*}


令 $2x = t$，则：


\begin{equation*}
\int_0^{\frac{\pi}{2}} \ln(\sin 2x)\,\mathrm{d}x = \frac{1}{2}\int_0^\pi \ln(\sin t)\,\mathrm{d}t = \frac{1}{2} \cdot 2\int_0^{\frac{\pi}{2}} \ln(\sin t)\,\mathrm{d}t = I
\end{equation*}


代回得 $2I = I - \frac{\pi}{2}\ln 2 \implies I = -\frac{\pi}{2}\ln 2$。

\end{thmbox}
\begin{thmbox}[模型[2]: 积分不等式深水区证明工具箱]

\textbf{[1]} 柯西-施瓦茨(Cauchy-Schwarz)积分不等式：
若 $f(x), g(x)$ 在 $[a, b]$ 上可积，则：


\begin{equation*}
\left(\int_a^b f(x)g(x)\,\mathrm{d}x\right)^2 \leqslant \left(\int_a^b f^2(x)\,\mathrm{d}x\right)\left(\int_a^b g^2(x)\,\mathrm{d}x\right)
\end{equation*}


等号成立充要条件: 存在常数 $k$ 使得 $f(x) = k g(x)$ 几乎处处成立。

典型应用: 设 $f(x) > 0$ 连续，证明 $\int_a^b f(x)\,\mathrm{d}x \int_a^b \frac{1}{f(x)}\,\mathrm{d}x \geqslant (b-a)^2$。

只需在柯西不等式中取 $g(x) = \frac{1}{\sqrt{f(x)}}$，代入立得 $(b-a)^2$。

\textbf{[2]} 变限积分微分方程构造法：
待证形如 $\int_0^x f(t)\,\mathrm{d}t \leqslant g(x)$ 时，构造辅助函数：


\begin{equation*}
\Phi(x) = g(x) - \int_0^x f(t)\,\mathrm{d}t
\end{equation*}


求导 $\Phi'(x) = g'(x) - f(x)$，通过一阶、二阶导数符号单调性由 $\Phi(0) = 0$ 推定 $\Phi(x) \geqslant 0$。

\end{thmbox}
\begin{thmbox}[模型[3]: 黎曼-勒贝格引理与高频振荡积分极限]

\textbf{[1]} 黎曼-勒贝格引理(Riemann-Lebesgue Lemma)：
若 $f(x)$ 在 $[a, b]$ 上黎曼可积，则当频率 $\lambda \to +\infty$ 时：


\begin{equation*}
\lim_{\lambda \to +\infty} \int_a^b f(x)\sin(\lambda x)\,\mathrm{d}x = 0, \qquad \lim_{\lambda \to +\infty} \int_a^b f(x)\cos(\lambda x)\,\mathrm{d}x = 0
\end{equation*}


数学本质: 高频振荡函数在极短区间内正负面积近乎完美抵消。

\textbf{[2]} 周期权重平均值极限定理：
若 $f(x)$ 在 $[a, b]$ 上连续，$g(x)$ 为以 $T$ 为周期的连续函数，则：


\begin{equation*}
\lim_{n \to \infty} \int_a^b f(x)g(nx)\,\mathrm{d}x = \left(\frac{1}{T}\int_0^T g(t)\,\mathrm{d}t\right) \int_a^b f(x)\,\mathrm{d}x
\end{equation*}


此结论是快速求解考研中乘积极限选择题的秘密大招。

\end{thmbox}

\section{3.变上限积分深水区方程与不可积函数极限}

\begin{thmbox}[模型[1]: 变限积分内不可积函数的泰勒级数展开内插法]

\textbf{[1]} 核心痛点与思维破局：
当遇到不可积函数(如 $e^{-t^2}, \frac{\sin t}{t}, \sqrt{1+t^3}, \ln(1+\sin t)$)作为变限积分被积函数，且洛必达求导后分子分母次数极高时，洛必达法则往往导致求导灾难。

破局心法: 泰勒公式直接在积分号内部展开，逐项积分脱积分号！

\textbf{[2]} 标准操作范例：
计算极限 $\lim_{x \to 0} \frac{\int_0^x \sin(t^2)\,\mathrm{d}t}{x^3}$：

在积分号内，当 $t \to 0$ 时，$\sin(t^2) = t^2 - \frac{t^6}{6} + o(t^6)$：


\begin{equation*}
\int_0^x \sin(t^2)\,\mathrm{d}t = \int_0^x \left(t^2 + o(t^2)\right)\,\mathrm{d}t = \frac{x^3}{3} + o(x^3)
\end{equation*}


瞬间得出极限值：


\begin{equation*}
= \lim_{x \to 0} \frac{\frac{x^3}{3} + o(x^3)}{x^3} = \frac{1}{3}
\end{equation*}


\end{thmbox}
\begin{thmbox}[模型[2]: 隐函数变上限积分高阶导数求解]

\textbf{[1]} 模型特征:
方程形式为 $y + \int_0^y e^{t^2}\,\mathrm{d}t = x + \sin x$，求 $\left.\frac{\mathrm{d}^2y}{\mathrm{d}x^2}\right|_{x=0}$。

\textbf{[2]} 规范求解SOP：
第一步: 确定初始点坐标，代入 $x=0$ 得 $y + 0 = 0 \implies y(0) = 0$。

第二步: 两边对 $x$ 求一阶导数：


\begin{equation*}
y' + e^{y^2} \cdot y' = 1 + \cos x \implies y'(1 + e^{y^2}) = 1 + \cos x
\end{equation*}


代入 $x=0, y=0$：$y'(0)(1+1) = 2 \implies y'(0) = 1$。

第三步: 两端对 $x$ 再次求导：


\begin{equation*}
y''(1+e^{y^2}) + y'(2y y' e^{y^2}) = -\sin x
\end{equation*}


代入 $x=0, y=0, y'(0)=1$：


\begin{equation*}
y''(0)(1+1) + 1 \cdot 0 = 0 \implies y''(0) = 0
\end{equation*}


\end{thmbox}

\section{4.反常积分计算深水区与含参变量积分}

\begin{thmbox}[模型[1]: 费曼积分号下求导法(参数积分法)实战]

\textbf{[1]} 狄利克雷积分推导：
求反常积分 $I = \int_0^{+\infty} \frac{\sin x}{x}\,\mathrm{d}x$。

构造含参辅助积分：引入指数收敛因子 $e^{-\alpha x}$ ($\alpha \geqslant 0$)：


\begin{equation*}
I(\alpha) = \int_0^{+\infty} e^{-\alpha x}\frac{\sin x}{x}\,\mathrm{d}x
\end{equation*}


对参数 $\alpha$ 求导并交换积分与导数次序：


\begin{equation*}
I'(\alpha) = \int_0^{+\infty} \frac{\partial}{\partial \alpha}\left(e^{-\alpha x}\right)\frac{\sin x}{x}\,\mathrm{d}x = -\int_0^{+\infty} e^{-\alpha x}\sin x\,\mathrm{d}x
\end{equation*}


此积分可用循环分部积分算出：


\begin{equation*}
I'(\alpha) = -\frac{1}{\alpha^2 + 1}
\end{equation*}


两边对 $\alpha$ 积分：


\begin{equation*}
I(\alpha) = -\arctan \alpha + C
\end{equation*}


确定常数: 当 $\alpha \to +\infty$ 时，$e^{-\alpha x} \to 0 \implies \lim_{\alpha \to +\infty} I(\alpha) = 0$。


\begin{equation*}
0 = -\frac{\pi}{2} + C \implies C = \frac{\pi}{2}
\end{equation*}


因此 $I(\alpha) = \frac{\pi}{2} - \arctan \alpha$。

令 $\alpha = 0$，即得狄利克雷积分极值：


\begin{equation*}
I = I(0) = \frac{\pi}{2}
\end{equation*}


\textbf{[2]} 高斯积分参数推广：
已知 $\int_0^{+\infty} e^{-x^2}\,\mathrm{d}x = \frac{\sqrt{\pi}}{2}$。

令 $I(a) = \int_0^{+\infty} e^{-a x^2}\,\mathrm{d}x = \frac{\sqrt{\pi}}{2\sqrt{a}}$。

两边对 $a$ 求 $n$ 阶导数，即可瞬时求出一切形如 $\int_0^{+\infty} x^{2n} e^{-a x^2}\,\mathrm{d}x$ 的超难积分。

\end{thmbox}
\begin{thmbox}[模型[2]: 瑕点与无穷端点双侧奇性的阶数双侧压制模型]

\textbf{[1]} 双侧奇点识别：
对于积分 $J = \int_0^{+\infty} \frac{x^p}{1+x^q}\,\mathrm{d}x$ ($q > 0$)：

下限 0 可能为瑕点，上限 $+\infty$ 为无穷限。必须在 $x=1$ 处拆分为两部分：


\begin{equation*}
J = \int_0^1 \frac{x^p}{1+x^q}\,\mathrm{d}x + \int_1^{+\infty} \frac{x^p}{1+x^q}\,\mathrm{d}x
\end{equation*}


\textbf{[2]} 双侧收敛充要条件判定：
\textbf{(1)}  在 $x \to 0^+$ 处：$1+x^q \sim 1$，被积函数等价于 $x^p = \frac{1}{x^{-p}}$。由瑕积分 $p$-审敛法，收敛充要条件为 $-p < 1 \implies p > -1$。

\textbf{(2)}  在 $x \to +\infty$ 处：$1+x^q \sim x^q$，被积函数等价于 $\frac{x^p}{x^q} = \frac{1}{x^{q-p}}$。由无穷积分 $p$-审敛法，收敛充要条件为 $q - p > 1$。

综上，$J$ 收敛的充要条件为：


\begin{equation*}
-1 < p < q - 1
\end{equation*}

\end{thmbox}

\end{document}
